\chapter{Considerações Finais}
\label{cap:consideracoes}

Este projeto propõe o desenvolvimento de uma solução escalável baseada em Realidade Virtual para treinamento em segurança do trabalho no setor da construção civil. Sua principal contribuição consiste em oferecer uma alternativa tecnológica capaz de simular situações de risco de forma controlada, imersiva e repetível, com potencial para reduzir significativamente a ocorrência de acidentes em canteiros de obras. Ao capacitar os trabalhadores por meio de experiências práticas em ambientes virtuais, busca-se não apenas ampliar o conhecimento técnico, mas também promover mudanças de comportamento em relação à identificação e prevenção de riscos ocupacionais.

Entre as limitações identificadas, destaca-se o custo inicial dos equipamentos de realidade virtual imersiva, que pode restringir sua adoção por parte de empresas de pequeno e médio porte. Além disso, a implementação de tecnologias imersivas em ambientes tradicionais de formação profissional ainda demanda um processo de adaptação cultural, exigindo tempo e estratégias pedagógicas apropriadas para garantir o engajamento dos usuários.

As perspectivas futuras para o aprimoramento da proposta incluem a integração com modelos BIM (\textit{Building Information Modeling}), de forma a permitir a criação de cenários personalizados baseados em projetos reais. Também se planeja incorporar sensores e rastreadores físicos --- como \textit{trackers} acoplados aos Equipamentos de Proteção Individual (EPIs) --- para aumentar o realismo tátil da simulação e expandir os pontos de coleta de dados avaliativos. Outra frente de desenvolvimento, ainda em estágio conceitual, contempla o uso de inteligência artificial para enriquecer a atuação de um instrutor virtual, que futuramente poderá ser capaz de interagir com os participantes durante a simulação, respondendo a dúvidas sobre normas técnicas e procedimentos de segurança. Embora essa funcionalidade não esteja prevista na versão atual do sistema, trata-se de um potencial evolutivo importante, especialmente para contextos educacionais mais autônomos. Além disso, tecnologias baseadas em IA poderão ser exploradas para gerar variações dinâmicas dos cenários de simulação, evitando que os participantes decorem os roteiros e promovendo uma experiência de aprendizado mais realista e desafiadora. Por fim, espera-se estabelecer parcerias com empresas e instituições reguladoras, a fim de validar a solução em contextos reais e fomentar sua adoção como ferramenta complementar nos programas de capacitação em segurança do trabalho.
